% ══════════════════════════════════════════════════════════════════════
% Gradient Atoms / Gradient SAE — Method and Results
% ══════════════════════════════════════════════════════════════════════

\subsection{EKFAC Eigenprojection}

For each training document $i$, we compute the cross-entropy loss gradient with respect to each tracked MLP weight matrix $W_m \in \mathbb{R}^{d_{\text{out}} \times d_{\text{in}}}$ (gate\_proj, up\_proj, down\_proj across $L{=}24$ layers, giving $M{=}72$ modules). The EKFAC decomposition~\citep{george2018fast} provides eigendecompositions of the Kronecker-factored Fisher:
\begin{equation}
    F_m \approx (Q_m^S \Lambda_m^S {Q_m^S}^\top) \otimes (Q_m^A \Lambda_m^A {Q_m^A}^\top)
\end{equation}
where $Q_m^A \in \mathbb{R}^{d_{\text{in}} \times d_{\text{in}}}$ and $Q_m^S \in \mathbb{R}^{d_{\text{out}} \times d_{\text{out}}}$ are activation and gradient eigenvectors, with corresponding eigenvalues $\Lambda_m^A$ and $\Lambda_m^S$.

For each module $m$, the Kronecker eigenvalues are $\lambda_{m,jk} = \Lambda^S_{m,j} \cdot \Lambda^A_{m,k}$, giving $d_{\text{out}} \times d_{\text{in}}$ eigenvalues per module. We retain the top-$k$ by magnitude ($k{=}50$ in our experiments, yielding $72 \times 50 = 3{,}600$ projected dimensions). The projection of document $i$'s gradient for module $m$ is:
\begin{equation}
    \tilde{g}_{i,m} = {Q_m^S}^\top \nabla_{W_m} \mathcal{L}_i \, Q_m^A \in \mathbb{R}^{d_{\text{out}} \times d_{\text{in}}}
\end{equation}
We flatten and select the top-$k$ components by Kronecker eigenvalue magnitude, concatenating across modules to form $\tilde{\mathbf{g}}_i \in \mathbb{R}^{Mk}$.

\paragraph{Preconditioning.} The standard EKFAC preconditioning scales each component by $1/\sqrt{|\lambda| + \epsilon}$, making the projected space approximately isotropic (equal curvature in all directions). This is optimal for influence functions~\citep{grosse2023studying} but, as we show, suboptimal for sparse decomposition: it amplifies noise in low-eigenvalue directions. We compare preconditioned and raw (unscaled) projections.

\paragraph{Optimised extraction.} Rather than computing the full eigenbasis projection $\tilde{g}_{i,m} \in \mathbb{R}^{d_{\text{out}} \times d_{\text{in}}}$ and selecting components, we precompute the unique eigenvector subsets needed for the top-$k$ Kronecker indices and project only onto those (${\sim}50$ columns each of $Q_m^A$ and $Q_m^S$), reducing per-module cost from $O(d_{\text{out}}^2 d_{\text{in}} + d_{\text{out}} d_{\text{in}}^2)$ to $O(k \cdot d_{\text{out}} \cdot d_{\text{in}})$.

\subsection{Normalisation}

We compare three normalisation strategies for the projected gradient vectors before sparse decomposition:

\paragraph{Unit normalisation.} $\hat{\mathbf{g}}_i = \tilde{\mathbf{g}}_i / \|\tilde{\mathbf{g}}_i\|_2$. Discards gradient magnitude; decomposes directions only.

\paragraph{Log normalisation.} $\hat{\mathbf{g}}_i = \frac{\tilde{\mathbf{g}}_i}{\|\tilde{\mathbf{g}}_i\|_2} \cdot \log(1 + \|\tilde{\mathbf{g}}_i\|_2)$. Preserves relative magnitudes on a compressed scale. Documents with large gradients (high training signal) receive proportionally more weight in the decomposition without outlier domination.

\paragraph{No normalisation.} $\hat{\mathbf{g}}_i = \tilde{\mathbf{g}}_i$. Preserves raw magnitudes. Unstable for dictionary learning (causes numerical issues in resampling) but viable for SAEs.

\subsection{Sparse Decomposition Methods}

\subsubsection{FISTA Dictionary Learning}

We seek a dictionary $\mathbf{D} \in \mathbb{R}^{K \times d}$ and sparse codes $\mathbf{Z} \in \mathbb{R}^{N \times K}$ minimising:
\begin{equation}
    \min_{\mathbf{D}, \mathbf{Z}} \frac{1}{N} \sum_{i=1}^{N} \left\| \hat{\mathbf{g}}_i - \sum_{j=1}^{K} z_{ij} \mathbf{d}_j \right\|_2^2 + \alpha \|z_i\|_1 \quad \text{s.t.} \quad \|\mathbf{d}_j\|_2 = 1 \;\forall j
\end{equation}
Sparse coding uses FISTA (Fast Iterative Shrinkage-Thresholding Algorithm) with soft thresholding. For scalability to $K{=}16{,}000$ atoms, we avoid forming the $K \times K$ matrix $\mathbf{D}\mathbf{D}^\top$ by computing gradients via the residual: $\nabla_{\mathbf{z}} = (\mathbf{Z}\mathbf{D} - \mathbf{X})\mathbf{D}^\top$, with cost $O(BKd)$ per iteration. The Lipschitz constant is estimated via power iteration on $\mathbf{D}\mathbf{D}^\top$.

\paragraph{Dead atom resampling.} Atoms activated fewer than a threshold number of times are reinitialised from training points with high reconstruction error, sampled proportional to $\|\hat{\mathbf{g}}_i - \hat{\mathbf{g}}_i^{\text{recon}}\|_2^2$.

\subsubsection{TopK Sparse Autoencoder}

Following \citet{gao2024scaling}, we train a TopK SAE with learned encoder:
\begin{align}
    \mathbf{z}_i &= \text{TopK}\!\left(\text{ReLU}\!\left(\mathbf{W}_{\text{enc}}(\hat{\mathbf{g}}_i - \mathbf{b}_{\text{dec}}) + \mathbf{b}_{\text{enc}}\right), K\right) \\
    \hat{\mathbf{g}}_i^{\text{recon}} &= \mathbf{W}_{\text{dec}} \mathbf{z}_i + \mathbf{b}_{\text{dec}}
\end{align}
where $\text{TopK}$ retains only the $K$ largest activations per input, setting the rest to zero. This fixes the sparsity level without requiring tuning of an L1 penalty $\alpha$.

\paragraph{AuxK dead feature recovery.} For features with activation count below a threshold (``dead'' features), we compute an auxiliary reconstruction of the residual $\mathbf{r}_i = \hat{\mathbf{g}}_i - \hat{\mathbf{g}}_i^{\text{recon}}$ using only the dead features:
\begin{equation}
    \mathcal{L}_{\text{aux}} = \frac{1}{N} \sum_{i=1}^{N} \left\| \mathbf{r}_i - \mathbf{W}_{\text{dec}} \, \text{TopK}_{\text{dead}}(\mathbf{W}_{\text{enc}} \mathbf{r}_i, K_{\text{aux}}) \right\|_2^2
\end{equation}
The total loss is $\mathcal{L} = \mathcal{L}_{\text{recon}} + \frac{1}{32} \mathcal{L}_{\text{aux}}$. Decoder columns are projected to unit norm after each optimisation step.

\subsection{Explained Variance}

We report reconstruction quality as explained variance:
\begin{equation}
    \text{EV} = 1 - \frac{\sum_{i=1}^{N} \|\hat{\mathbf{g}}_i - \hat{\mathbf{g}}_i^{\text{recon}}\|_2^2}{\sum_{i=1}^{N} \|\hat{\mathbf{g}}_i\|_2^2}
\end{equation}
where $\hat{\mathbf{g}}_i^{\text{recon}} = \mathbf{D}^\top \mathbf{z}_i$ (FISTA) or $\mathbf{W}_{\text{dec}} \mathbf{z}_i + \mathbf{b}_{\text{dec}}$ (SAE).

\subsection{Coherence}

For each atom/feature $j$, we measure coherence over \emph{all} activating documents $\mathcal{S}_j = \{i : z_{ij} \neq 0\}$ using the efficient mean-of-normed estimator:
\begin{equation}
    \text{coherence}(j) = \frac{\left\| \frac{1}{|\mathcal{S}_j|} \sum_{i \in \mathcal{S}_j} \frac{\tilde{\mathbf{g}}_i}{\|\tilde{\mathbf{g}}_i\|} \right\|^2 - \frac{1}{|\mathcal{S}_j|}}{\frac{|\mathcal{S}_j| - 1}{|\mathcal{S}_j|}}
\end{equation}
which equals the mean pairwise cosine similarity and is computable in $O(|\mathcal{S}_j| \cdot d)$ rather than $O(|\mathcal{S}_j|^2 \cdot d)$.

% ══════════════════════════════════════════════════════════════════════
% Results Table
% ══════════════════════════════════════════════════════════════════════

\subsection{Results}

\begin{table}[t]
\centering
\caption{Explained variance of sparse gradient decompositions on SmolLM2-1.7B-Instruct (1.04M SmolTalk documents, 72 MLP modules, 8$\times$ H200 GPUs).}
\label{tab:gradient-atoms-sweep}
\small
\begin{tabular}{llccccc}
\toprule
\textbf{Method} & \textbf{Features} & \textbf{Dims} & \textbf{Precond.} & \textbf{Norm} & \textbf{EV (\%)} & \textbf{Time} \\
\midrule
\multicolumn{7}{l}{\textit{FISTA Dictionary Learning}} \\
\midrule
FISTA ($\alpha{=}0.1$)   & 500      & 3{,}600  & \cmark & unit & 28.2 & 4m \\
FISTA ($\alpha{=}0.1$)   & 16{,}000 & 3{,}600  & \cmark & unit & 39.3 & 25m \\
FISTA ($\alpha{=}0.05$)  & 16{,}000 & 3{,}600  & \cmark & unit & 2.0  & 25m \\
FISTA ($\alpha{=}0.1$)   & 2{,}000  & 3{,}600  & \cmark & log  & 53.8 & 22m \\
FISTA ($\alpha{=}0.07$)  & 16{,}000 & 3{,}600  & \cmark & unit & 38.4 & 33m \\
\midrule
\multicolumn{7}{l}{\textit{TopK Sparse Autoencoder ($K{=}64$, AuxK $K_{\text{aux}}{=}256$)}} \\
\midrule
TopK SAE   & 16{,}384 & 3{,}600  & \cmark & log  & 62.9 & 3m \\
TopK SAE   & 16{,}384 & 14{,}400 & \cmark & log  & 46.4 & 11m \\
TopK SAE   & 65{,}536 & 14{,}400 & \cmark & log  & 51.8 & 43m \\
\rowcolor{green!10}
TopK SAE   & 16{,}384 & 3{,}600  & \xmark & log  & \textbf{72.3} & \textbf{50s} \\
\bottomrule
\end{tabular}
\vspace{0.5em}

\footnotesize
\textbf{Key findings:} (1)~TopK SAE outperforms FISTA by $+$24pp at matched feature count (62.9 vs 39.3\%). (2)~Log-normalisation (preserving gradient magnitudes) improves over unit normalisation by $+$15pp. (3)~Removing EKFAC preconditioning improves EV by $+$9pp by concentrating signal in high-eigenvalue directions (357$\times$ dynamic range). (4)~Increasing EKFAC dimensions from 3{,}600 to 14{,}400 \emph{decreases} EV at matched overcompleteness---additional low-eigenvalue components add noise amplified by preconditioning. (5)~Sparsity penalty $\alpha$ is critical for FISTA: $\alpha{=}0.05$ collapses to 2\% EV (codes activate all atoms with negligible coefficients).
\end{table}

% ══════════════════════════════════════════════════════════════════════
% Preconditioning vs Raw: Interpretability Analysis
% ══════════════════════════════════════════════════════════════════════

\subsection{Effect of EKFAC Preconditioning on Feature Interpretability}

While removing EKFAC preconditioning improves explained variance (72.3\% vs 64.0\%), we find it significantly reduces feature diversity and interpretability. We compare the two settings using four quantitative metrics on TopK SAE features (16{,}384 features, $K{=}64$, log-normalised inputs).

\begin{table}[t]
\centering
\caption{Interpretability comparison: preconditioned (isotropic) vs unpreconditioned (raw Fisher) gradient space. Both use identical TopK SAE architecture (16{,}384 features, $K{=}64$, log-norm).}
\label{tab:preconditioning-interpretability}
\small
\begin{tabular}{lcc}
\toprule
\textbf{Metric} & \textbf{Preconditioned} & \textbf{Unpreconditioned} \\
\midrule
Explained variance (\%)          & 64.0  & \textbf{72.3} \\
Active features ($>$100 docs)    & \textbf{11{,}035} & 1{,}488 \\
Features with coherence $>$0.5   & 32    & \textbf{138} \\
Mean coherence (active)          & 0.069 & \textbf{0.200} \\
\midrule
\multicolumn{3}{l}{\textit{Diversity metrics (top 200 features by coherence)}} \\
\midrule
Unique top-3 keyword combos      & \textbf{71\%} & 66\% \\
Top keyword concentration        & \textbf{42\%} & 48\% \\
Mean inter-feature cosine sim    & \textbf{0.0004} & 0.0376 \\
Feature pairs with $\cos > 0.5$  & \textbf{0} & 4 \\
Normalised activation entropy    & 0.80  & \textbf{0.93} \\
\bottomrule
\end{tabular}
\end{table}

\paragraph{Inter-feature similarity.} The most striking difference is in decoder vector similarity: preconditioned features exhibit $94\times$ lower mean cosine similarity (0.0004 vs 0.038) with zero redundant pairs ($\cos > 0.5$). This confirms that preconditioning produces features capturing genuinely distinct behavioral modes, while the unpreconditioned space yields many near-duplicate features aligned with the dominant Fisher eigendirections.

\paragraph{Feature utilisation.} Preconditioning yields $7.4\times$ more active features (11{,}035 vs 1{,}488 with $>$100 activating documents). In the unpreconditioned space, the 357$\times$ dynamic range between eigenvalue-scaled dimensions causes the SAE encoder to concentrate activations in a small subset of features. The isotropic space distributes activations more democratically.

\paragraph{Coherence vs diversity tradeoff.} Unpreconditioned features achieve higher per-feature coherence (mean 0.20 vs 0.07) because they capture the dominant gradient directions where many documents agree. However, this comes at the cost of diversity: the top keyword appears in 48\% of features vs 42\%, and only 66\% of top-3 keyword combinations are unique vs 71\%.

\paragraph{Interpretation.} The EKFAC preconditioning transforms gradient space from a \emph{variance-dominated basis} to a \emph{behavioural basis}. Two gradients close in preconditioned space induce similar changes in model predictions, regardless of raw magnitude. This makes preconditioned features semantically meaningful by construction. We recommend using preconditioned features for interpretability analysis and unpreconditioned features when maximum reconstruction fidelity is required (e.g., for steering vector extraction).

\paragraph{Practical recommendation.} An effective pipeline decouples the two: train the SAE on preconditioned gradients for feature discovery and interpretability, then at steering time, unproject the discovered features through the raw (unpreconditioned) Fisher eigenbasis to maximise the behavioural impact of each steering vector.
